\section{结构、数据与控制冒险}

冒险表示下一拍若照常推进，就会使用不存在、过期或错误路径上的信息。

\begin{longtable}{L{0.18\linewidth}L{0.36\linewidth}L{0.36\linewidth}}
\toprule
类型 & 产生原因 & 典型处理 \\
\midrule
结构冒险 & 两个阶段同时争同一硬件端口 & 增加端口、分离资源或停顿 \\
数据冒险 & 消费者早于生产者提交结果 & 转发、互锁停顿、编译器调度 \\
控制冒险 & 下一 PC 尚未确定 & 预测执行；错误时冲刷并重定向 \\
\bottomrule
\end{longtable}

考虑相邻两条指令：

\begin{lstlisting}
ADD R3, R1, R2
SUB R5, R3, R4
\end{lstlisting}

ADD 的结果在 EX 末端已经算出，但要到 WB 才写入寄存器堆。SUB 在下一拍进入 EX 时若只读 ID 阶段取得的旧 R3 就会出错。EX/MEM 到 ALU 输入的旁路可以把结果直接送给 SUB；旁路选择条件必须同时比较目标寄存器号、源寄存器号和“前条指令确实写寄存器”三项。

load-use 更难：

\begin{lstlisting}
LOAD R3, [R1]
ADD  R5, R3, R4
\end{lstlisting}

加载数据通常到 MEM 末端才出现，紧随其后的 ADD 在同一拍开头已经需要输入，即使有旁路也来不及，必须插入至少一个气泡。若数据 Cache miss，固定的一拍停顿还要升级为“保持请求状态直到响应真正返回”。

\section{精确异常要求按程序顺序提交}

流水线里可能同时有多条指令处于不同阶段。发生异常时，架构状态必须看起来像：异常指令之前的都已完成，异常指令及之后的都没有产生可见副作用。实现上应先阻止年轻指令写寄存器、写存储器或发出不可撤销 MMIO，再保存异常 PC 和原因，最后跳到异常入口。

乱序核心进一步把“执行完成”和“按序提交”分开。重排序缓冲记录程序顺序，执行单元可以先算就绪的独立操作，但只有队首指令确认无异常时才能把结果变成架构可见状态。store buffer 也必须等到提交条件成立，才允许写操作越过不可恢复的边界。

\begin{keyidea}
性能优化可以改变内部执行顺序，不能改变程序看到的提交顺序、异常位置和设备副作用顺序。
\end{keyidea}
